% !TEX root = ../../../main.tex

\subsection{拼接句子视频整体识别结果}

在词级分类模型测试的基础上，
本文进一步对拼接句子视频识别流程进行整体测试。

该测试用于验证系统在连续短句场景下的端到端识别能力，
包括样本拼接、滑动窗口识别、词段合并、候选序列生成以及受约束语义修正等环节。

本节采用 core61 测试集作为拼接句子视频测试样例。

该测试集共包含 61 条核心短句，
短句长度覆盖 2 至 6 个 Gloss 不等，
能够反映系统在不同长度连续动作序列下的识别表现。

需要说明的是，
由于公开可用的单词级手语视频资源有限，
当前实验采用的是 25 类小词表模型。

词表规模和单词样本数量会直接影响句子级识别效果。

因此，
本节的重点并不是将拼接句子识别结果理解为大规模通用手语识别能力，
而是验证在当前小词表和有限样本条件下，
系统是否能够跑通从连续视频到 Gloss 序列输出的完整链路，
以及候选序列生成和语义修正是否具有进一步提升空间。

为了避免单次随机划分或单次运行结果对实验结论造成偏倚，
本文不以某一个随机种子下的结果作为主结论，
而是统计 10 组随机种子的平均结果。

对于每条测试样例，
系统首先基于词级分类模型输出原始识别序列 \texttt{rawSequence}。

随后，
系统根据每个词段的 Top-K 候选结果生成候选 Gloss 序列集合。

最后，
DeepSeek 在候选集合中选择一个候选编号，
后端再根据该候选编号得到修正后的 Gloss 序列。

表~\ref{tab:chapter06-spliced-sentence-multi-seed-result} 给出了 10 组随机种子下的平均测试结果。

\begin{longtable}{L{0.32\textwidth} L{0.22\textwidth} L{0.36\textwidth}}
  \caption{拼接句子视频多随机种子平均结果表}
  \label{tab:chapter06-spliced-sentence-multi-seed-result}\\
  \toprule
  指标 & 平均结果 & 说明 \\
  \midrule
  total & 61 & 每组随机种子下的测试样例总数 \\
  raw\_exact & 35.3 / 61 & 原始检测序列平均完全匹配数量 \\
  raw\_accuracy & 57.87\% & 原始检测序列平均完全匹配率 \\
  candidate\_oracle\_exact & 47.1 / 61 & 候选集合中平均包含正确序列的样例数 \\
  candidate\_oracle\_accuracy & 77.21\% & 候选集合平均理论完全匹配率 \\
  selected\_exact & 36.8 / 61 & DeepSeek 候选重排序后平均完全匹配数量 \\
  selected\_accuracy & 60.33\% & DeepSeek 候选重排序后平均完全匹配率 \\
  net\_gain & +1.5 & DeepSeek 候选重排序带来的平均净提升样例数 \\
  \bottomrule
\end{longtable}

由表~\ref{tab:chapter06-spliced-sentence-multi-seed-result} 可知，
在 10 组随机种子的平均结果中，
原始拼接句子视频识别的完全匹配率为 57.87\%。

这说明在当前小词表条件下，
系统已经能够在一部分连续短句样例中生成完全正确的 Gloss 序列。

同时，
该结果也说明原始识别仍然受到训练资源规模、动作边界、局部误判和词段后处理策略的影响。

DeepSeek 候选重排序后的平均完全匹配率为 60.33\%，
相比原始识别结果有一定提升。

候选集合理论完全匹配率为 77.21\%，
明显高于原始识别结果和 DeepSeek 实际选择结果。

这说明，
虽然原始 Top-1 序列中存在一定错误，
但系统通过保留词段 Top-K 候选并组合生成候选序列，
能够在相当一部分错误样例中保留正确答案。

因此，
候选序列生成模块本身是有效的，
它证明了当前识别链路中存在较大的可恢复空间。

在这个意义上，
候选集合理论上限并不是最终系统效果，
而是用于说明：
只要候选排序策略进一步优化，
句子级识别结果仍有继续提升的可能。

图~\ref{fig:chapter06-core61-average-summary} 进一步展示了 core61 句子集在多随机种子下的平均完全匹配准确率。

\WhuAutoFigure{0.82\textwidth}{0.42\textheight}{figures/chapter06/core61-average-summary.png}{Core61 句子集平均识别准确率汇总}{fig:chapter06-core61-average-summary}

由图~\ref{fig:chapter06-core61-average-summary} 可以看出，
候选上限明显高于原始识别和 DeepSeek 候选重排序结果。

这表明当前候选序列生成方法能够有效扩大正确答案空间。

同时，
DeepSeek 候选重排序结果高于原始识别结果，
说明受约束语义修正策略具有一定的稳定增益。

但是，
DeepSeek 实际选择结果与候选上限之间仍存在差距。

该差距并不意味着语义修正机制无效，
而是说明在当前小词表、小样本和候选信息有限的条件下，
候选排序仍然是系统后续优化的重要方向。

图~\ref{fig:chapter06-core61-accuracy-comparison} 展示了不同随机种子下原始识别、DeepSeek 候选重排序和候选上限三类结果的对比。

\WhuAutoFigure{0.92\textwidth}{0.48\textheight}{figures/chapter06/core61-accuracy-comparison.png}{Core61 句子集多随机种子识别准确率对比}{fig:chapter06-core61-accuracy-comparison}

从图~\ref{fig:chapter06-core61-accuracy-comparison} 可以看出，
不同随机种子下三类结果的整体趋势基本一致。

候选上限始终明显高于原始识别结果，
说明 Top-K 候选保留和候选序列生成策略具有稳定价值。

DeepSeek 候选重排序结果整体略高于原始识别结果，
说明受约束语义修正能够在不同随机种子下保持一定正向作用。

本文前期也尝试过更激进的语义修正策略。

但是，
在当前 25 类小词表和有限训练样本条件下，
激进策略容易过度依赖语言合理性，
从而产生脱离视觉证据的改写结果。

这种情况下，
修正后的序列反而可能低于原始识别序列。

因此，
本文最终采用候选集合约束下的保守重排序策略。

该策略的目标不是让大语言模型自由改写识别结果，
而是在尽量保持视觉证据忠实性的前提下，
从有限候选中选择更合理的序列。

从连续 10 组随机种子的结果看，
该策略能够带来稳定但有限的平均提升，
说明 DeepSeek 语义修正在当前系统中具有一定潜力。

结合实验结果可以归纳出，
拼接句子视频识别中主要存在以下几类错误现象。

表~\ref{tab:chapter06-sentence-error-types} 给出了句子视频识别错误现象归纳。

\begin{longtable}{L{0.22\textwidth} L{0.34\textwidth} L{0.36\textwidth}}
  \caption{句子视频识别错误现象归纳表}
  \label{tab:chapter06-sentence-error-types}\\
  \toprule
  错误类型 & 典型表现 & 可能原因 \\
  \midrule
  插入误报 & 序列中多出无关 Gloss，如两个动作之间出现额外词段 & 动作边界或过渡帧被误识别为具体词汇 \\
  漏词 & 期望序列中的某个 Gloss 未出现在检测序列中 & 词段置信度不足、动作持续时间短或被 NMS 抑制 \\
  重复或冗余输出 & 同一语义片段附近出现重复词或冗余词 & 滑动窗口重叠导致同一动作被多个窗口重复响应 \\
  局部替换错误 & 某个 Gloss 被识别为动作相近或语义无关的其他词 & 训练样本不足、手型相近或动作轨迹相似，导致词级分类混淆 \\
  候选排序未命中 & 正确序列存在于候选集合中，但最终未被 DeepSeek 选中 & 候选排序阶段对视觉证据、置信度信息和语义合理性的综合判断仍不充分 \\
  \bottomrule
\end{longtable}

由表~\ref{tab:chapter06-sentence-error-types} 可知，
插入误报和冗余输出是拼接句子视频识别中的重要问题。

这与滑动窗口识别机制的特点有关。

为了覆盖连续动作中的词汇边界，
系统需要使用相互重叠的窗口进行密集预测。

这种方式提高了动作片段被捕获的概率，
但也会使词与词之间的过渡帧形成额外候选窗口，
从而引入边界插入词或冗余输出。

漏词和局部替换错误则更多反映了词级分类模型和后处理策略的局限。

由于高质量单词级视频资源较少，
当前词级模型的训练样本规模有限。

当某个词段持续时间较短、
关键点检测质量不足，
或对应类别与其他动作类别相似时，
该词段可能被过滤、被抑制，
或被识别为其他 Gloss。

候选排序未命中则说明，
候选集合中存在正确答案并不等于最终系统一定能够选中正确答案。

在当前策略下，
DeepSeek 被限制在程序生成的候选集合中进行选择，
这能够避免自由生成带来的不受控改写，
但也要求候选信息本身能够充分表达视觉证据和局部置信度。

如果候选描述不足，
或原始序列在自然语言层面仍然可以解释，
DeepSeek 可能倾向于保留较保守的候选。

总体来看，
拼接句子视频整体测试表明，
当前系统已经能够完成从连续视频到 Gloss 序列的端到端识别。

候选集合的理论命中率显著高于原始识别完全匹配率，
说明 Top-K 候选保留和候选序列生成具有继续优化价值。

DeepSeek 候选重排序带来了稳定但有限的平均提升，
说明受约束语义修正在当前小词表条件下具有一定潜力。

同时，
实验结果也说明当前系统的主要限制来自两个方面。

一方面，
高质量单词级手语视频资源不足，
导致词表规模较小，
词级模型对相似动作的区分能力仍有提升空间。

另一方面，
候选排序阶段尚未充分利用候选集合中已经存在的正确序列。

后续若要进一步提升句子级识别效果，
需要继续扩充高质量训练资源，
提高词级模型的分类能力，
并在保持候选约束的前提下改进语义修正模块对视觉证据和语言合理性的综合判断。